\section{Sensores e Calibração}

Os principais sensores estão resumidos na Tabela \ref{tab:sensores_ambientais}.
Um nó pode integrar diferentes sensores, desde que cada transdutor seja instalado de acordo com a variável medida e com suas condições de exposição. Sensores de temperatura do ar requerem abrigo ventilado e proteção contra radiação solar direta; sondas de solo exigem contato estável com o substrato; e sensores de qualidade da água requerem imersão adequada e manutenção periódica para controle da bioincrustação \cite{5597912,10.1145/3005719,s23146332}.

Quanto à qualidade metrológica, sensores de baixo custo permitem aumentar a densidade espacial do monitoramento. Sua equivalência a instrumentos de referência depende de calibração, co-localização e validação nas condições de operação \cite{atmos10090506,KANG2022151769}.
Em particular, estudos sobre sensores de baixo custo para qualidade do ar identificam a temperatura, a umidade, a deriva temporal, as sensibilidades cruzadas e o modelo de calibração como fatores que influenciam a resposta dos dispositivos \cite{atmos10090506,KANG2022151769}.
Com base nesses fatores, recomenda-se co-localizar uma parcela dos nós com um instrumento de referência, documentar o modelo e a versão da calibração aplicada e repetir a verificação durante a implantação para identificar alterações de desempenho \cite{atmos10090506,KANG2022151769}.

\begin{table*}[!t]
\centering
\caption{Variáveis, princípios de medição e requisitos de implantação para o monitoramento ambiental.}
\label{tab:sensores_ambientais}
\footnotesize
\setlength{\tabcolsep}{4pt}
\renewcommand{\arraystretch}{1.18}

\begin{tabularx}{\textwidth}{
    @{}
    L{2.0cm}
    L{3.6cm}
    Y
    Y
    @{}
}
\toprule
\textbf{Grupo} &
\textbf{Variáveis monitoradas} &
\textbf{Princípios de medição e sensores} &
\textbf{Requisitos de implantação e manutenção} \\
\midrule

Água
\cite{10.1145/3005719} &
Temperatura, pH, turbidez, oxigênio dissolvido e condutividade elétrica &
Termistor ou RTD; eletrodo de pH; turbidímetro óptico; sensor óptico ou eletroquímico de oxigênio dissolvido; célula de condutividade &
Controle de bioincrustação; compensação de temperatura; limpeza periódica; verificação e recalibração dos sensores \\

\addlinespace

Eventos naturais
\cite{10.1145/1134680.1134685} &
Nível da água, vibração do solo, pressão acústica, fumaça e velocidade do vento &
Transdutor de pressão ou sensor ultrassônico; acelerômetro ou geofone; microfone; detector óptico de fumaça; anemômetro &
Definição e validação dos limiares; sincronização temporal; aquisição disparada por evento; armazenamento temporário local \\

\addlinespace

Microclima
\cite{HART2006177,5597912,10.1145/570738.570751} &
Temperatura do ar, umidade relativa e pressão atmosférica &
Termistor, RTD ou sensor digital de temperatura; sensor capacitivo de umidade; barômetro MEMS &
Abrigo meteorológico ventilado; altura e exposição padronizadas; proteção contra radiação solar direta e condensação \\

\addlinespace

Qualidade do ar
\cite{atmos10090506,KANG2022151769} &
\(\mathrm{PM}_{2{,}5}\), \(\mathrm{PM}_{10}\), \(\mathrm{CO}_{2}\), \(\mathrm{CO}\), \(\mathrm{NO}_{2}\), \(\mathrm{O}_{3}\) e compostos orgânicos voláteis &
Contador óptico de partículas; sensor NDIR; células eletroquímicas; sensores de óxido metálico &
Co-localização com instrumento de referência; calibração; compensação de temperatura e umidade; avaliação de deriva e sensibilidade cruzada \\

\addlinespace

Radiação solar e precipitação
\cite{5597912} &
Iluminância, irradiância solar e precipitação &
Fotodiodo ou luxímetro; piranômetro; pluviômetro de báscula &
Nivelamento; campo de visão desobstruído; controle de sombreamento; limpeza e inspeção periódicas \\

\addlinespace

Solo
\cite{5597912,s23146332} &
Teor volumétrico de água, temperatura e condutividade elétrica &
Sensor capacitivo, FDR ou TDR; termistor ou RTD; eletrodos de condutividade &
Contato adequado com o solo; calibração para o tipo de solo; avaliação do efeito da salinidade; distribuição espacial representativa dos pontos \\

\bottomrule
\end{tabularx}

\vspace{3pt}
\parbox{\textwidth}{\footnotesize
RTD: detector resistivo de temperatura;
MEMS: sistemas microeletromecânicos;
NDIR: infravermelho não dispersivo;
FDR: reflectometria no domínio da frequência;
TDR: reflectometria no domínio do tempo.
}
\end{table*}